\begin{resumen}
	El problema de la satisfacibilidad booleana (SAT) es un problema NP-completo que consiste en determinar si existe una asignación de valores de verdad a las variables de una fórmula booleana que haga que la fórmula sea verdadera. La rama de la teoría de lenguajes formales es una rama de la ciencia de la computación que estudia los lenguajes formales y sus propiedades. El objetivo de este trabajo es vincular el problema de la satisfacibilidad booleana con la teoría de lenguajes formales para construir el lenguaje de las interpretaciones que satisfacen una fórmula. 

	En este trabajo, se propone un enfoque que consiste en construir un autómata finito que reconoce el lenguaje de las interpretaciones que satisfacen a cada clausula independientemente, asumiendo que cada cláusula puede satisfacerse con una interpretación diferente. Luego, se intersecta dicho autómata con un lenguaje que obliga a que cada cláusula se satisfaga con la misma interpretación. Esto permite que el problema de la satisfacibilidad booleana pueda resolverse determinando si el lenguaje resultante es vacío o no. 

	A partir de esta construcción se demuestra que, para cualquier formalismo capaz de generar este lenguaje de repeticiones y cerrado bajo intersección con lenguajes regulares, al menos uno de los siguientes problemas es NP-duro: construir la intersección con un lenguaje regular o decidir el vacío del lenguaje resultante.

	Posteriormente se estudia el caso de las gramáticas libres de contexto múltiples paralelas (pMCFG), un formalismo reportado en la literatura como cerrado bajo intersección con lenguajes regulares y con problema de vacío decidible en tiempo polinomial. Se muestra que el procedimiento de intersección propuesto en la literatura para este formalismo es incorrecto mediante un contraejemplo, y se demuestra que el problema de determinar la intersección entre una pMCFG y un lenguaje regular es NP-duro.
\end{resumen}

\begin{abstract}
	The Boolean satisfiability problem (SAT) is an NP-complete problem that consists of determining whether there exists an assignment of truth values to the variables of a Boolean formula such that the formula evaluates to true. Formal language theory is a branch of computer science that studies formal languages and their properties. The objective of this work is to relate the Boolean satisfiability problem to formal language theory in order to construct the language of the interpretations that satisfy a formula.

	In this work, an approach is proposed that consists of constructing a finite automaton that recognizes the language of interpretations that satisfy each clause independently, assuming that each clause may be satisfied by a different interpretation. Then, this automaton is intersected with a language that enforces all clauses to be satisfied by the same interpretation. This allows the Boolean satisfiability problem to be solved by determining whether the resulting language is empty or not.

	From this construction, it is shown that for any formalism capable of generating this repetition language and closed under intersection with regular languages, at least one of the following problems is NP-hard: constructing the intersection with a regular language or deciding emptiness of the resulting language.

	Subsequently, the case of parallel multiple context-free grammars (pMCFGs) is studied. This formalism has been reported in the literature as being closed under intersection with regular languages and having a polynomial-time emptiness problem. It is shown that the intersection procedure proposed in the literature for this formalism is incorrect by means of a counterexample, and it is proved that determining the intersection between a pMCFG and a regular language is NP-hard.

\end{abstract}